\newpage
\begin{center}
  \textbf{\large ПРИЛОЖЕНИЕ А ИЗОБРАЖЕНИЕ КАК НАБОР ТОЧЕК}
\end{center}
\refstepcounter{chapter}
\addcontentsline{toc}{chapter}{ПРИЛОЖЕНИЕ А ИЗОБРАЖЕНИЕ КАК НАБОР ТОЧЕК}

\label{sec:set_of_points}
Важную роль в рассматриваемом в работе методе играет особый способ представления изображений
как набора точек на основе контекстной кластеризации~\cite{ma2023image}. Идея заключается
в использовании нового подхода для экстракции признаков при работе с визуальными данными,
позволяющем увеличить эффективность при качестве сравнимом с широко применяемыми
сверточными нейронными сетями~\cite{simonyan2014very} и визуальными трансформерами~\cite{dosovitskiy2020image}.

Конвенционально изображения представлялись в качестве прямоугольного $h\times w$ трехканального массива
$I \in \mathbb{R}^{h\times w \times 3}$. В данном подходе предлагается использовать пространственную
информацию о положении пикселя в качестве непосредственно первичного признака, а все изображение
рассматривать как неупорядоченный набор точек с обобщенными признаками $I^P \in \mathbb{R}^{n_p \times 5}$,
где $n_p = h\cdot w$. При этом пространственная информация каждого пикселя $I_{i,j}$
представляется следующим образом: $\begin{bmatrix} \frac{i}{h}-0.5 & \frac{j}{w}-0.5 \end{bmatrix}$.

После преобразования изображений производится иерархическая экстракция признаков. Первым шагом
является редуцирование набора точек посредством равномерного выбора ключевых точек и
проецирования некоторой окрестности при помощи линейного слоя в единое представление.
После редуцирования применяется процедура контекстной кластеризации, что и является ключевой
идеей предлагаемого метода. Редуцированные точки проецируются в пространство сравнения при помощи
полносвязной нейронной сети. Среди полученных точек $I^p_s \in \mathbb{R}^{n_r \times d}$
(количество признаков может меняться при проецировании) также равномерно выбираются
центры будущих $c$ кластеров --- их признаки рассчитываются как среднее среди $k$ ближайших соседей.
Затем рассчитывается карта сходства $S \in \mathbb{R}^{c\times n_r}$ на основе
косинусного расстояния между центрами кластеров и точками. Таким образом,
используя карту сходства можем определить принадлежность точек к кластерам.
Имея кластер, состоящий из $m$ точек и соответствующие им значения сходства $s \in \mathbb{R}^m$
производится проецирование в пространство значений $I^p_v \in \mathbb{R}^{n_r\times d'}$ (также проецируется центр кластера $v_c$).
Рассчитывается агрегированный признак класса:
\begin{equation*}
  g = \frac{1}{C} \left( v_c + \sum_{i=1}^m \sigma(a s_i + \beta) \cdot v_i \right), \quad  C = 1 + \sum_{i=1}^m \sigma(a s_i + \beta),
\end{equation*}
где $v_i$ --- значение для $i$-й точки кластера, $\sigma$ --- сигмоидальная функция, $\alpha$ и $\beta$ --- обучаемые параметры.
После получения $g$ предлагается обновить признаки каждой точки кластера. Пусть $p_i \in \mathbb{R}^{d}$ ---
исходный вектор признаков $i$-й точки. Итоговое обновление выполняется формулой~\eqref{eq:coc_update}.
\begin{equation}
  \label{eq:coc_update}
  p'_i \;=\; p_i \;+\;
  \mathrm{FC}\Bigl(\sigma(\alpha\, s_i + \beta)\,\cdot\, g\Bigr),
\end{equation}
где $\mathrm{FC}$ --- полносвязный слой, отображающий из $\mathbb{R}^{d'}$ в $\mathbb{R}^{d}$.
Таким образом, каждая точка получает <<распространённый>> признак, пропорциональный величине сходства $s_i$.
Данная процедура повторяется несколько раз в зависимости от количества блоков в иерархическом
модуле.